\section{Model and Preliminaries}
\label{sec:model}
Our model consists of $N$ model developers (e.g., OpenAI, Anthropic, DeepMind), each of which competes on $K$ dimensions (e.g., math, coding, science). 
Developer $i$ has a production function $F_{i}: \mathbb{R}_{+}^{K} \rightarrow \mathbb{R}_{+}^{K}$, and investment $\mathbf{x}_i$, which produces capabilities $\bthetai := \Fi(\bxi)$. 
Investment at level $\bxi$ results in a cost $c(\onenorm \bxi)$, where $c$ is a strictly increasing, differentiable, continuous, and unbounded cost function with $c(0) = 0$ (per-developer cost functions $c_i$ can be folded into the model without changing the equilibria; see Appendix~\ref{app:CommonCost_proof}). 
We model each developer as maximizing its expected win rate.
We use expected weighted pairwise win rate as a simplified objective for relative competitive performance, rather than assuming that developers literally optimize a particular public leaderboard.
This objective directly corresponds to maximizing weighted category-level pairwise win probabilities under Bradley–Terry/Elo-style evaluation (see Appendix~\ref{app:elo}). 
Here, win-rate is weighted according to $\bw \geq 0, \sum_{k} w_{k} = 1$.
We define the utility as 
\begin{equation}
\label{eq:utility}
u_{i}(\mathbf{x}) = \frac{1}{N-1} \sum_{i' \neq i} \sum_{k} w_{k} \sigma(\Fi(\bxi)_k - F_{i'}(\bx_{i'})_k)- c(\sum_{k} x_{i,k}).
\end{equation}
Here, $\sigma$ denotes the win-probability according to some pairwise preference model. 
We assume $\sigma$ is strictly increasing, twice differentiable, bounded in $[0,1]$, $\sigma'(x)>0$, and symmetric where $\sigma(x)+\sigma(-x)=1$ (e.g., Bradley-Terry).
Our model abstracts from the details of developers' internal R\&D processes and instead captures three features of model development: developers care about their capabilities relative to competitors, allocate costly investment across multiple capabilities, and differ in how effectively investment translates into capability improvements. 
We do not require that developers directly optimize a particular leaderboard; tournament-style evaluation provides a tractable representation of competition based on relative model performance. 
The model can also be viewed as a multidimensional contest, related to Colonel Blotto games~\citep{colonel_blotto}, in which investment is transformed into capabilities through heterogeneous production functions \(F_i\).

Let $\mathbf{x} = (\mathbf{x}_1,\ldots,\mathbf{x}_{N})$ be a Nash equilibrium, and let $\Hij(\mathbf{x}) = \cos(\bthetai,\btheta_{i'})$ be the pairwise monoculture for developers $i$ and $i'$, with $\Hij(\mathbf{x})= 1$ indicating that $\bthetai$ and $\btheta_{i'}$ are parallel. 
We can view $\btheta_i$ as consisting of a direction $\mathbf{d}^*_i \in S_{+}^{K-1}$ and a magnitude $\alpha^*_{i}$, with only the direction impacting $\Hij$. 
We focus on cosine similarity here to capture capability profiles rather than capability levels; while two developers might differ in their overall level of output, they can still be similar if their relative category outputs are similar~\citep{monoculture_outcome}. 
We return to similarity in the overall level of output in Section~\ref{sec:dynamics}. 
We define overall monoculture as $H(\mathbf{x}) = \frac{1}{{N \choose 2}} \sum_{i<i'} \Hij(\mathbf{x})$, with $H=1$ indicating that all developers have parallel $\bthetai$. 
When exactly one capability vector is zero we define $H_{i,i'}(\mathbf{x}^*)=0$, and when both are zero we define $H_{i,i'}(\mathbf{x}^*)=1$.